% Part: turing-machines
% Chapter: machines-computations
% Section: well-behaved-machines

\documentclass[../../../include/open-logic-section]{subfiles}

\begin{document}

\olfileid[pt]{tur}{mac}{dis}
\olsection{Máquinas Disciplinadas}

\begin{explain}
Na \olref[hal]{sec}, consideramos máquinas de Turing que têm um
único estado de parada designado~$h$---tais máquinas são garantidas de
parar, se pararem de fato, no estado~$h$. Dessa forma, máquinas com um
único estado de parada são mais ``disciplinadas'' do que permitimos que
máquinas de Turing sejam em geral. Há outras restrições que podemos
impor ao comportamento das máquinas de Turing. Por exemplo, também não
proibimos máquinas de Turing de apagar o marcador de fim de fita
no quadrado~$0$, ou de tentar mover-se à esquerda a partir do quadrado~$0$. (Nossa
definição afirma que a cabeça simplesmente permanece no quadrado~$0$ nesse
caso; outras definições fazem a máquina parar.) Da mesma forma, às vezes é
desejável poder supor que uma máquina de Turing, se
parar de fato, para no quadrado~$1$.
\end{explain}

\begin{defn}\ollabel{defn:disciplined}
Uma máquina de Turing~$M$ é \emph{disciplinada} sse
\begin{enumerate}
    \item ela tem um único estado de parada designado~$h$,
    \item ela para, se parar de fato, enquanto examina o quadrado~$1$,
    \item ela nunca apaga o símbolo $\TMendtape$ no quadrado~$0$, e
    \item ela nunca tenta mover-se à esquerda a partir do quadrado~$0$.
\end{enumerate}
\end{defn}

\begin{explain}
Já discutimos que qualquer máquina de Turing pode ser transformada em
uma com o mesmo comportamento, mas com um estado de parada designado. Isso é
feito simplesmente adicionando um novo estado~$h$, e adicionando uma instrução
$\delta(q, \sigma) = \tuple{h, \sigma, N}$ para qualquer par
$\tuple{q,\sigma}$ onde o $\delta$ original é indefinido. É
verdade, embora tedioso de provar, que qualquer máquina de Turing~$M$ pode ser
transformada em uma máquina de Turing disciplinada~$M'$ que para nas mesmas
entradas e produz a mesma saída. Por exemplo, se a máquina de Turing
para e não está no quadrado~$1$, podemos adicionar algumas instruções
para fazer a cabeça mover-se à esquerda até encontrar o marcador de fim de fita, depois
mover um quadrado à direita, e então parar. Deixamos a você pensar
sobre como as outras condições podem ser tratadas.
\end{explain}

\begin{ex}
    Em \olref{fig:adder-disc}, transformamos a máquina de adição
    de \olref[una]{ex:adder} em uma máquina disciplinada.
    \begin{figure}\[
\begin{tikzpicture}[->,>=stealth',shorten >=1pt,auto,node distance=2.8cm,
                    semithick]
  \tikzstyle{every state}=[fill=none,draw=black,text=black]

  \node[initial,state]         (A)              {$q_0$};
  \node[state]         (B) [right of=A] {$q_1$};
  \node[state]         (C) [below right of=B] {$q_2$};
  \node[state]         (D) [below left of=C] {$q_3$};
  \node[state]         (H) [left of=D] {$h$};

  \path (A) edge node {\TMtrans{\TMblank}{\TMstroke}{\TMstay}} (B)
           edge [loop below] node {\TMtrans{\TMstroke}{\TMstroke}{\TMright}} (B)
        (B) edge [loop below] node {\TMtrans{\TMstroke}{\TMstroke}{\TMright}} (B)
            edge node {\TMtrans{\TMblank}{\TMblank}{\TMleft}} (C)
        (C) edge node {\TMtrans{\TMstroke}{\TMblank}{\TMleft}} (D)
        (D) edge [loop above] node {\TMtrans{\TMstroke}{\TMstroke}{\TMleft}} (D)
        edge node {\TMtrans{\TMendtape}{\TMendtape}{\TMright}} (H);
\end{tikzpicture}
\]\caption{Uma máquina de adição disciplinada}
\ollabel{fig:adder-disc}
\end{figure}
\end{ex}

\begin{prop}\ollabel{prop:disciplined} Para toda máquina de Turing~$M$,
há uma máquina de Turing disciplinada~$M'$ que para com saída~$O$
se $M$~para com saída~$O$, e não para se $M$~não para.
Em particular, qualquer função $f\colon\Nat^n \to \Nat$ computável por uma
máquina de Turing também é computável por uma máquina de Turing disciplinada.
\end{prop}

\begin{prob}\label{tur:mac:dis:prob:disc-succ}
Dê uma máquina disciplinada que computa $f(x) = x+1$.
\end{prob}


\begin{prob}\label{tur:mac:dis:prob:copier}
Encontre uma máquina disciplinada que, quando iniciada com entrada $\TMstroke^n$
produz saída $\TMstroke^n \concat \TMblank \concat \TMstroke^n$.
\end{prob}

\end{document}
